\documentclass{beamer}
% \usepackage{moreverb} 
\usepackage{listings}
\usepackage{mflogo}

\mode<presentation> {
  \usetheme{Warsaw}
  \setbeamercovered{opaque}
}

\usebackgroundtemplate{\includegraphics[width=\paperwidth]{format/libresoft-bg.png}}
\usepackage[spanish]{babel}
\usepackage[utf8]{inputenc}
\usepackage{graphics}
\usepackage{amssymb} % Simbolos matematicos

%% Metadatos del PDF.
\hypersetup{  
  pdftitle={KVM: Curso práctico},
  pdfauthor={Jose Castro},
  pdfcreator={GSyC/Libresoft},
  pdfproducer=PDFLaTeX,
  pdfsubject={Master on Free Software},
}
%%

\defbeamertemplate*{footline}{shadow theme}
{%
  \leavevmode%
  \hbox{\begin{beamercolorbox}[wd=.5\paperwidth,ht=2.5ex,dp=1.125ex,leftskip=.1cm plus1fil,rightskip=.2cm]{author in head/foot}%
\includegraphics[scale=0.40]{format/cc-by-80x15.png} \hspace{0.05cm}
	Miguel Vidal / Jose Castro
  \end{beamercolorbox}%
  \begin{beamercolorbox}[wd=.5\paperwidth,ht=2.5ex,dp=1.125ex,leftskip=.3cm,rightskip=.3cm plus1fil]{title in head/foot}%
    \usebeamerfont{title in head/foot}\insertshorttitle%
    \hspace{2cm} \usebeamerfont{author in head/foot}\insertframenumber\,/\,\inserttotalframenumber\hfill
  \end{beamercolorbox}}%
  \vskip0pt%
}


\begin{document}

\title{KVM: Curso práctico}
\subtitle{System Integration}
% \institute{\texttt{http://gsyc.urjc.es/\~{}mvidal} \\ Twitter: \texttt{@mvidallopez}}
\author{Miguel Vidal, Jose Castro} 
\date{\footnotesize{Master on Free Software \\ April 13rd, 2012}}
\author{Miguel Vidal \hspace{1cm} Jose Castro \\
\hspace{0.5mm} {\tiny Twitter: @mvidallopez \hspace{1.1cm}Twitter: @jfcastroluis}
}


\frame{
\maketitle
\begin{center}
\includegraphics[width=6cm]{format/gsyc-urjc}
\end{center}
}

%% License slide
\begin{frame}
  \vspace{2cm}
  \begin{flushright}
    {\small \copyright{} 2010-2012 Miguel Vidal, Jose Castro} \\
%    \vspace{0.25cm}
    \medskip
    {\scriptsize This work is licensed under \\ a Creative Commons Attribution 3.0 License}
%    \vspace{0.10cm}
  \end{flushright}
  \begin{flushright}
    \href{http://creativecommons.org/licenses/by/3.0/es}{\includegraphics[width=2cm]{format/cc-by.png}} \\
    {\tiny \url{http://creativecommons.org/licenses/by/3.0}}
  \end{flushright}
\end{frame}%%

\usebackgroundtemplate{}

\AtBeginSection[]
{
\begin{frame}<beamer>
\begin{center}
{\huge \insertsection}
\end{center}
\end{frame}
}

\AtBeginSection[]
{
  \begin{frame}<beamer>{Índice}
    \tableofcontents[currentsection]
  \end{frame}
}


\AtBeginSubsection[]
{
  \begin{frame}<beamer>{Índice}
    \tableofcontents[currentsection,currentsubsection]
  \end{frame}
}

%%


%%%%%%%%%%%%%%%%%%%%%%%%%%%%%%%%%%%%%%%%%%%%%%%%%%%%%%%%%%%%%%%%%%%%%%%

%%%%%%%%%%%%%%%%%%%%%%%%%%%%%%%%%%%%%%%%%%%%%%%%%%%%%%%%%%%%%%%%%%%%%%%
\section{Prerrequisitos}
%%%%%%%%%%%%%%%%%%%%%%%%%%%%%%%%%%%%%%%%%%%%%%%%%%%%%%%%%%%%%%%%%%%%%%%

\begin{frame}
  \frametitle{Instrucciones de virtualización}
  \begin{itemize}
    \item Activar las instrucciones en la BIOS
  \end{itemize}
  \bigskip
  \begin{block}{Comprobación}
    \texttt{\$ egrep -c '(vmx|svm)' /proc/cpuinfo}
  \end{block}
  \bigskip
  \begin{itemize}
    \item Con kernel Xen comprobar flag \texttt{hvm} en \texttt{/sys/hypervisor/properties/capabilities}
  \end{itemize}
\end{frame}
%%%%%%%%%%%%%%%%%%%%%%%%%%%%%%%%%%%%%%%%%%%%%%%%%%%%%%%%%%%%%%%%%%%%%%%

\begin{frame}
  \frametitle{64 bits}
  Se debe usar un kernel de 64 bits para:
  \begin{itemize}
    \item configurar más de 2GB a una máquina virtual
    \item poder crear máquinas virtuales de 32 y 64 bits
  \end{itemize}
  \begin{block}{Comprobación procesador}
    \texttt{\$ egrep -c 'lm' /proc/cpuinfo}
  \end{block}
  \begin{block}{Comprobación kernel}
    \texttt{\$ uname -m}
    \begin{itemize}
      \item i686 -- 32 bits
      \item x86\_64 -- 64 bits
    \end{itemize}
  \end{block}
\end{frame}
%%%%%%%%%%%%%%%%%%%%%%%%%%%%%%%%%%%%%%%%%%%%%%%%%%%%%%%%%%%%%%%%%%%%%%%

%%%%%%%%%%%%%%%%%%%%%%%%%%%%%%%%%%%%%%%%%%%%%%%%%%%%%%%%%%%%%%%%%%%%%%%
\section{Instalación}
%%%%%%%%%%%%%%%%%%%%%%%%%%%%%%%%%%%%%%%%%%%%%%%%%%%%%%%%%%%%%%%%%%%%%%%

\begin{frame}
  \frametitle{Instalación}
  \begin{block}{Instalación del sistema}
    \texttt{\# apt-get install qemu-kvm|kvm libvirt-bin [bridge-utils]}
    \begin{itemize}
      \item qemu-kvm --- módulo kvm.ko
      \item libvirt-bin --- gestión de la virtualización
      \item bridge-utils --- utilidades de red para \textit{bridges}
    \end{itemize}
  \end{block}
  \bigskip
  Los usuarios que vayan a gestionar las máquinas virtuales han de pertenecer a los grupos \texttt{libvirtd} y \texttt{kvm}
\end{frame}
%%%%%%%%%%%%%%%%%%%%%%%%%%%%%%%%%%%%%%%%%%%%%%%%%%%%%%%%%%%%%%%%%%%%%%%

\begin{frame}[fragile]
  \frametitle{Comprobación}
  \begin{block}{virsh}
    \texttt{\$ virsh -c qemu:///system list} \\
    \small{\begin{verbatim}
Id Name                  State
----------------------------------
    \end{verbatim}}
  \end{block}
  \bigskip
  \begin{block}{kvm-ok}
    \texttt{\$ kvm-ok} \\
    \small{\begin{verbatim}
INFO: /dev/kvm exists
KVM acceleration can be used
    \end{verbatim}}
  \end{block}
\end{frame}
%%%%%%%%%%%%%%%%%%%%%%%%%%%%%%%%%%%%%%%%%%%%%%%%%%%%%%%%%%%%%%%%%%%%%%%

\begin{frame}
  \frametitle{Herramientas}
  \begin{block}{virtinst}
    \texttt{\# apt-get install virtinst}
  \end{block}
  Este paquete provee las siguientes herramientas:
  \begin{itemize}
    \item \texttt{virt-install}
    \item \texttt{virt-clone}
    \item \texttt{virt-image}
    \item \texttt{virt-convert}
  \end{itemize}
\end{frame}
%%%%%%%%%%%%%%%%%%%%%%%%%%%%%%%%%%%%%%%%%%%%%%%%%%%%%%%%%%%%%%%%%%%%%%%

%%%%%%%%%%%%%%%%%%%%%%%%%%%%%%%%%%%%%%%%%%%%%%%%%%%%%%%%%%%%%%%%%%%%%%%
\section{Configuración}
%%%%%%%%%%%%%%%%%%%%%%%%%%%%%%%%%%%%%%%%%%%%%%%%%%%%%%%%%%%%%%%%%%%%%%%

\begin{frame}[fragile]
  \frametitle{Pools almacenamiento}
  \begin{block}{Pool default}
    \texttt{\$ virsh pool-start default} \\
    \texttt{\$ virsh pool-autostart default} \\
    \medskip
    \texttt{\$ virsh pool-list $--$all} \\
    \small{\begin{verbatim}
Name                 State      Autostart 
----------------------------------------- 
default              active     yes
    \end{verbatim}}
  \end{block}
  \begin{center}
    \texttt{/etc/libvirt/storage/default.xml}
  \end{center}
\end{frame}
%%%%%%%%%%%%%%%%%%%%%%%%%%%%%%%%%%%%%%%%%%%%%%%%%%%%%%%%%%%%%%%%%%%%%%%

\begin{frame}[fragile]
  \frametitle{Red}
  \begin{block}{Red interna}
    \texttt{\$ virsh net-start default} \\
    \texttt{\$ virsh net-autostart default} \\
    \medskip
    \texttt{\$ virsh net-list $--$all} \\
    \small{\begin{verbatim}
Name                 State      Autostart
-----------------------------------------
default              active     yes
    \end{verbatim}}
  \end{block}
  \bigskip
  \begin{itemize}
    \item Las máquinas virtuales reciben una IP por DHCP
    \item El anfitrión hace NAT para la conexión de los invitados
  \end{itemize}
  \begin{center}
    \texttt{/etc/libvirt/qemu/networks/default.xml}
  \end{center}
\end{frame}
%%%%%%%%%%%%%%%%%%%%%%%%%%%%%%%%%%%%%%%%%%%%%%%%%%%%%%%%%%%%%%%%%%%%%%%

%%%%%%%%%%%%%%%%%%%%%%%%%%%%%%%%%%%%%%%%%%%%%%%%%%%%%%%%%%%%%%%%%%%%%%%
\section{Gestión de Máquinas Virtuales}
%%%%%%%%%%%%%%%%%%%%%%%%%%%%%%%%%%%%%%%%%%%%%%%%%%%%%%%%%%%%%%%%%%%%%%%

\begin{frame}
  \frametitle{Creación}
  \begin{block}{Creación VM}
    \scriptsize{\texttt{\$ virt-install $--$connect qemu:///system \textbackslash \\
    $--$name debian $--$ram 512 $--$os-type linux \textbackslash \\
    $--$os-variant debiansqueeze \textbackslash \\
    $--$disk path=/var/lib/libvirt/images/debian.img,size=5 \textbackslash \\
    $--$cdrom /home/jfcastro/isos/debian-6.0.4-amd64-netinst.iso}}
  \end{block}
  \begin{itemize}
    \item Configuración: \texttt{/etc/libvirt/qemu/debian.xml}
    \item Imagen: \texttt{/var/lib/libvirt/images.debian.img}
  \end{itemize}
\end{frame}
%%%%%%%%%%%%%%%%%%%%%%%%%%%%%%%%%%%%%%%%%%%%%%%%%%%%%%%%%%%%%%%%%%%%%%%

\begin{frame}[fragile]
  \frametitle{Arranque}
  \begin{block}{Arranque VM}
    \texttt{\$ virsh -c qemu:///system start debian} \\
    \small{\begin{verbatim}
Domain debian started
    \end{verbatim}}
  \end{block}
  \bigskip
  \begin{block}{Arranque automático}
    \texttt{\$ virsh -c qemu:///system autostart debian} \\
    \small{\begin{verbatim}
Domain debianbase marked as autostarted
    \end{verbatim}}
  \end{block}
  Crea un enlace simbólico en \texttt{/etc/libvirt/qemu/autostart}
\end{frame}
%%%%%%%%%%%%%%%%%%%%%%%%%%%%%%%%%%%%%%%%%%%%%%%%%%%%%%%%%%%%%%%%%%%%%%%

\begin{frame}
  \frametitle{Consola y conexiones}
  \begin{block}{Instalación}
    \texttt{\# apt-get install virt-viewer}
  \end{block}
  \begin{block}{Consola}
    \texttt{\$ virt-viewer -c qemu:///system debian}
  \end{block}
  \bigskip
  \begin{block}{Conexión local}
    \texttt{\$ virsh -c qemu:///system}
  \end{block}
  \begin{block}{Conexión remota}
    \texttt{\$ virsh -c qemu+ssh://user@domain:port/system}
  \end{block}
\end{frame}
%%%%%%%%%%%%%%%%%%%%%%%%%%%%%%%%%%%%%%%%%%%%%%%%%%%%%%%%%%%%%%%%%%%%%%%

\begin{frame}
  \frametitle{Parada}
  \begin{block}{Parada VM}
    \texttt{\$ virsh -c qemu:///system shutdown debian}
  \end{block}
  \bigskip
  \begin{block}{Botonazo VM}
    \texttt{\$ virsh -c qemu:///system destroy debian}
  \end{block}
\end{frame}
%%%%%%%%%%%%%%%%%%%%%%%%%%%%%%%%%%%%%%%%%%%%%%%%%%%%%%%%%%%%%%%%%%%%%%%

\begin{frame}
  \frametitle{Borrado}
  \begin{block}{Borrado VM}
    \texttt{\$ virsh -c qemu:///system destroy debian} \\
    \texttt{\$ virsh -c qemu:///system undefine debian} \\
  \end{block}
  \bigskip
  \begin{block}{Borrado ficheros}
    \texttt{\# rm -r /var/lib/libvirt/images/debian.img} \\
    \texttt{\# rm /etc/libvirt/qemu/autostart/debian.xml} \\
    \texttt{\# rm /etc/libvirt/qemu/debian.xml} \\
  \end{block}
\end{frame}
%%%%%%%%%%%%%%%%%%%%%%%%%%%%%%%%%%%%%%%%%%%%%%%%%%%%%%%%%%%%%%%%%%%%%%%

\begin{frame}
  \frametitle{Comandos básicos}
  \begin{itemize}
    \item \textbf{virsh help}: muestra la ayuda del comando
    \item \textbf{virsh list \-\-all}: listado de dominios activos
    \item \textbf{virsh dominfo}: muestra información de un dominio
    \item \textbf{virsh dumpxml}: muestra la configuración xml
    \item \textbf{virsh suspend}: parar temporalmente un dominio
    \item \textbf{virsh resume}: vuelve a la actividad un dominio suspendido
    \item \textbf{virsh save}: guarda un dominio ejecutando en un fichero
    \item \textbf{virsh restore}: crea un dominio desde un fichero
    \item \textbf{virsh migrate}: migra un dominio a otro host
  \end{itemize}
\end{frame}
%%%%%%%%%%%%%%%%%%%%%%%%%%%%%%%%%%%%%%%%%%%%%%%%%%%%%%%%%%%%%%%%%%%%%%%

\begin{frame}
  \frametitle{GUI}
  \begin{block}{virt-manager}
    \texttt{\# apt-get install virt-manager}
  \end{block}
  \bigskip
  \begin{itemize}
    \item Se puede gestionar pools, redes y máquinas virtuales
    \item Permite conexiones locales y remotas
    \item Es una interfaz gráfica realmente completa
  \end{itemize}
\end{frame}
%%%%%%%%%%%%%%%%%%%%%%%%%%%%%%%%%%%%%%%%%%%%%%%%%%%%%%%%%%%%%%%%%%%%%%%

%%%%%%%%%%%%%%%%%%%%%%%%%%%%%%%%%%%%%%%%%%%%%%%%%%%%%%%%%%%%%%%%%%%%%%%
\section{Almacenamiento externo}
%%%%%%%%%%%%%%%%%%%%%%%%%%%%%%%%%%%%%%%%%%%%%%%%%%%%%%%%%%%%%%%%%%%%%%%

\begin{frame}
  \frametitle{Creación fichero}
  \begin{block}{Creación fichero almacenamiento}
    \small{\texttt{\# dd if=/dev/zero of=/var/lib/libvirt/images/data.img} \textbackslash  \\
    \texttt{bs=1024 count=102400}}
  \end{block}
  \bigskip
  \begin{center}
    El fichero debe estar en uno de los pools de almacenamiento
  \end{center}
\end{frame}
%%%%%%%%%%%%%%%%%%%%%%%%%%%%%%%%%%%%%%%%%%%%%%%%%%%%%%%%%%%%%%%%%%%%%%%

\begin{frame}[fragile]
  \frametitle{Configuración máquina virtual}
  \begin{block}{Configuración almacenamiento}
    \texttt{\$ virsh -c qemu:///system shutdown debian} \\
    \texttt{\$ virsh -c qemu:///system edit debian} \\
    \small{\begin{verbatim}
<disk type='file' device='disk'>
   <driver name='qemu' type='raw'/>
   <source file='/var/lib/libvirt/images/data.img'/>
   <target dev='vdb' bus='virtio'/>
<disk>
   \end{verbatim}}
   \texttt{\$ virsh -c qemu:///system start debian}
  \end{block}
\end{frame}
%%%%%%%%%%%%%%%%%%%%%%%%%%%%%%%%%%%%%%%%%%%%%%%%%%%%%%%%%%%%%%%%%%%%%%%

\begin{frame}[fragile]
  \frametitle{Configuración dispositivo}
  \begin{block}{Creación partición}
    \texttt{\# cfdisk /dev/vdb}
  \end{block}
  \begin{block}{Formato partición}
    \texttt{\# mkfs.ext3 /dev/vdb1}
  \end{block}
  \begin{block}{Montaje sistema ficheros}
    \texttt{\# mount /dev/vdb1 /mnt} \\
    \texttt{\# vim /etc/fstab} \\
    \small{\begin{verbatim}
/dev/vdb1	/mnt	ext3	defaults	0	0
    \end{verbatim}}
  \end{block}
\end{frame}

\end{document}
